\begin{table}[htbp!]
    \centering
    \small
    \begin{tabular}{|l|r|r|r|}
    \hline
    \textbf{Line Item} & \textbf{Per Unit} & \textbf{70-Yr Couple} & \textbf{Family of 4} \\ \hline
    \multicolumn{4}{|c|}{\textit{Assumptions}} \\ \hline
    Household Composition & --- & 2 Adults (65+) & 2 Adults, 2 Children \\ \hline
    Assessed Home Value & --- & \$350,000 & \$350,000 \\ \hline
    \multicolumn{4}{|c|}{\textit{Revenue Generated}} \\ \hline
    Property Tax (2\% effective) & --- & \$7,000 & \$7,000 \\ \hline
    Sales Tax \& Other & \$1,996/capita & \$3,992 & \$7,983 \\ \hline
    \textbf{Total Revenue} & --- & \textbf{\$10,992} & \textbf{\$14,983} \\ \hline
    \multicolumn{4}{|c|}{\textit{Marginal Cost of Services}} \\ \hline
    Public Education (K--12) & \$6,860/pupil & \$0 & \$13,719 \\ \hline
    Other Local Services & \$6,555/household & \$6,555 & \$6,555 \\ \hline
    \textbf{Total Marginal Cost} & --- & \textbf{\$6,555} & \textbf{\$20,274} \\ \hline
    \multicolumn{4}{|c|}{\textit{Net Fiscal Impact}} \\ \hline
    \textbf{Surplus / (Deficit)} & --- & \textbf{+\$4,437} & \textbf{(\$5,291)} \\ \hline
    \end{tabular}
    \caption{Illustrative marginal fiscal impact of new residents on local government. \textit{Revenue}: Property tax assumes 2\% effective rate on assessed value. Sales tax \& other includes local sales taxes, charges, and fees (\$1,996/capita = \$662B total non-property local own-source revenue $\div$ 331.9M population). \textit{Expenditure}: Education calculated as local direct expenditure on elementary and secondary education (\$753B) divided by public K--12 enrollment (49.4M students). Other local services calculated as local direct general expenditure excluding education (\$1123B) divided by 129.9M households, covering police, fire, roads, sewerage, parks, courts, and administration---services that primarily serve households rather than individuals. Education expenditure reflects local-funded share only (45\%, per NCES). For other services, 55\% of K--12 spending is allocated to state/federal transfers (\$414B), leaving \$353B for non-education---implying a 76\% local share. Source: U.S. Census Bureau, 2021 Annual Survey of State and Local Government Finances; population and households from Census Bureau; enrollment from NCES.}
    \label{tab:fiscal_impact}
\end{table}